% GENERATED FILE. Edit canonical lessons, then run npm run generate:books.
% canonical-source-hash: fnv1a64:911b795592694bb2
% canonical-lessons: JA-C372-ekimae, JA-C372-baiten, JA-C372-machiawaseru, JA-C372-machiaishitsu, JA-C372-jidouhanbaiki

\chapter{In Front of the Station}
\label{ch:ja-in-front-of-the-station}
\hlchaptermodality{372}

\begin{chapteropening}
\textbf{By the end of this chapter:} \emph{I can arrange to meet someone at the station.}

Complete the last lesson of chapter 372: I can arrange to meet someone at the station.
\end{chapteropening}

\section[\texorpdfstring{\ja{えきまえ} (ekimae)}{えきまえ (ekimae)}]{\ja{えきまえ} (ekimae) — the area in front of a station}

\label{lesson:JA-C372-ekimae}



\begin{quote}
Before the new one: say the Japanese for serious, then the Japanese for selfish.
\end{quote}

\subsection*{You'll want to know: \ja{えきまえ}}
\textbf{\ja{えきまえ}} — \emph{ekimae} — \textquotedblleft{}the area in front of a station\textquotedblright{}.

\begin{quote}
\textbf{\ja{えきまえに} \ja{ゆうびんきょくが} \ja{あります。}} — \emph{ekimae ni yūbinkyoku ga arimasu} — \textquotedblleft{}There is a post office in front of the station.\textquotedblright{}
\end{quote}

\begin{cousinweb}
\textbf{\ja{えきまえ}} (\emph{ekimae}) is built from both layers of the vocabulary: the Sino-Japanese \emph{eki}, \textquotedblleft{}station\textquotedblright{} (Middle Chinese \emph{yek}), an on'yomi borrowed from Chinese, and the native Japanese (wago) \emph{mae}, \textquotedblleft{}front\textquotedblright{}, inherited rather than borrowed (JMdict; KANJIDIC2; Middle Chinese in Baxter's transcription).
\end{cousinweb}

\begin{grammarlens}[title={What you've built}]
One more word for getting around.
\end{grammarlens}

\subsection*{Your turn}
\begin{itemize}
\raggedright
  \item \emph{Say it:} \emph{ekimae}
  \item \emph{Say it:} \emph{ekimae}, once more
  \item \emph{Say it:} \emph{ekimae ni yūbinkyoku ga arimasu} — There is a post office in front of the station.
  \item \emph{Recall:} say the Japanese for serious, then the Japanese for selfish, then say \emph{ekimae} again
\end{itemize}

\begin{tcolorbox}[breakable,colback=teal!4,colframe=teal!35!black,title={Before you move on}]
What does \ja{えきまえ} mean? (\textquotedblleft{}The area in front of a station\textquotedblright{}.) Say it once more.
\end{tcolorbox}
\section[\texorpdfstring{\ja{ばいてん} (baiten)}{ばいてん (baiten)}]{\ja{ばいてん} (baiten) — a kiosk; a stall}

\label{lesson:JA-C372-baiten}



\begin{quote}
Before the new one: say the Japanese for selfish, then the Japanese for in front of a station.
\end{quote}

\subsection*{You'll want to know: \ja{ばいてん}}
\textbf{\ja{ばいてん}} — \emph{baiten} — \textquotedblleft{}a kiosk; a stall\textquotedblright{}.

\begin{quote}
\textbf{\ja{えきの} \ja{ばいてんで} \ja{しんぶんを} \ja{かいました。}} — \emph{eki no baiten de shinbun o kaimashita} — \textquotedblleft{}I bought a newspaper at the station kiosk.\textquotedblright{}
\end{quote}

\begin{cousinweb}
\textbf{\ja{ばいてん}} (\emph{baiten}) is Sino-Japanese (kango): two Chinese characters read with their on'yomi, readings borrowed from Chinese. They are \emph{bai}, \textquotedblleft{}sell\textquotedblright{} (Middle Chinese \emph{meaH}) and \emph{ten}, \textquotedblleft{}store\textquotedblright{} (\emph{temH}) (KANJIDIC2; Middle Chinese in Baxter's transcription).
\end{cousinweb}

\begin{grammarlens}[title={What you've built}]
One more word for getting around.
\end{grammarlens}

\subsection*{Your turn}
\begin{itemize}
\raggedright
  \item \emph{Say it:} \emph{baiten}
  \item \emph{Say it:} \emph{baiten}, once more
  \item \emph{Say it:} \emph{eki no baiten de shinbun o kaimashita} — I bought a newspaper at the station kiosk.
  \item \emph{Recall:} say the Japanese for selfish, then the Japanese for in front of a station, then say \emph{baiten} again
\end{itemize}

\begin{tcolorbox}[breakable,colback=teal!4,colframe=teal!35!black,title={Before you move on}]
What does \ja{ばいてん} mean? (\textquotedblleft{}A kiosk; a stall\textquotedblright{}.) Say it once more.
\end{tcolorbox}
\section[\texorpdfstring{\ja{まちあわせる} (machiawaseru)}{まちあわせる (machiawaseru)}]{\ja{まちあわせる} (machiawaseru) — to meet up (at an agreed place)}

\label{lesson:JA-C372-machiawaseru}



\begin{quote}
Before the new one: say the Japanese for in front of a station, then the Japanese for a kiosk.
\end{quote}

\subsection*{You'll want to know: \ja{まちあわせる}}
\textbf{\ja{まちあわせる}} — \emph{machiawaseru} — \textquotedblleft{}to meet up (at an agreed place)\textquotedblright{}.

\begin{quote}
\textbf{\ja{ともだちと} \ja{えきまえで} \ja{まちあわせました。}} — \emph{tomodachi to ekimae de machiawasemashita} — \textquotedblleft{}I met up with my friend in front of the station.\textquotedblright{}
\end{quote}

Polite \textbf{\ja{まちあわせます}}, \emph{machiawasemasu}; an ichidan verb. Past (polite): \textbf{\ja{ちちと} \ja{ははが} \ja{えきで} \ja{まちあわせました。}}, \emph{chichi to haha ga eki de machiawasemashita}.

\begin{cousinweb}
\textbf{\ja{まちあわせる}} (\emph{machiawaseru}) is a native Japanese word (wago), inherited rather than borrowed from Chinese, built from \emph{machi}, \textquotedblleft{}waiting\textquotedblright{} (from \emph{matsu}) and \emph{awaseru}, \textquotedblleft{}to bring together\textquotedblright{}; its kanji were attached later as kun'yomi (JMdict).
\end{cousinweb}

\begin{grammarlens}[title={What you've built}]
One more word for getting around.
\end{grammarlens}

\subsection*{Your turn}
\begin{itemize}
\raggedright
  \item \emph{Say it:} \emph{machiawaseru}
  \item \emph{Say it:} \emph{machiawaseru}, once more
  \item \emph{Say it:} \emph{tomodachi to ekimae de machiawasemashita} — I met up with my friend in front of the station.
  \item \emph{Recall:} say the Japanese for in front of a station, then the Japanese for a kiosk, then say \emph{machiawaseru} again
\end{itemize}

\begin{tcolorbox}[breakable,colback=teal!4,colframe=teal!35!black,title={Before you move on}]
What does \ja{まちあわせる} mean? (\textquotedblleft{}To meet up (at an agreed place)\textquotedblright{}.) Say it once more.
\end{tcolorbox}
\section[\texorpdfstring{\ja{まちあいしつ} (machiaishitsu)}{まちあいしつ (machiaishitsu)}]{\ja{まちあいしつ} (machiaishitsu) — a waiting room}

\label{lesson:JA-C372-machiaishitsu}



\begin{quote}
Before the new one: say the Japanese for a kiosk, then the Japanese for to meet up.
\end{quote}

\subsection*{You'll want to know: \ja{まちあいしつ}}
\textbf{\ja{まちあいしつ}} — \emph{machiaishitsu} — \textquotedblleft{}a waiting room\textquotedblright{}.

\begin{quote}
\textbf{\ja{まちあいしつで} \ja{でんしゃを} \ja{まちました。}} — \emph{machiaishitsu de densha o machimashita} — \textquotedblleft{}I waited for the train in the waiting room.\textquotedblright{}
\end{quote}

\begin{cousinweb}
\textbf{\ja{まちあいしつ}} (\emph{machiaishitsu}) is built from both layers of the vocabulary: the native Japanese (wago) \emph{machi}, \textquotedblleft{}waiting\textquotedblright{} (from \emph{matsu}), inherited rather than borrowed; the native Japanese (wago) \emph{ai}, \textquotedblleft{}meeting\textquotedblright{} (from \emph{au}), inherited rather than borrowed; and the Sino-Japanese \emph{shitsu}, \textquotedblleft{}room\textquotedblright{} (Middle Chinese \emph{syit}), an on'yomi borrowed from Chinese (JMdict; KANJIDIC2; Middle Chinese in Baxter's transcription).
\end{cousinweb}

\begin{grammarlens}[title={What you've built}]
One more word for getting around.
\end{grammarlens}

\subsection*{Your turn}
\begin{itemize}
\raggedright
  \item \emph{Say it:} \emph{machiaishitsu}
  \item \emph{Say it:} \emph{machiaishitsu}, once more
  \item \emph{Say it:} \emph{machiaishitsu de densha o machimashita} — I waited for the train in the waiting room.
  \item \emph{Recall:} say the Japanese for a kiosk, then the Japanese for to meet up, then say \emph{machiaishitsu} again
\end{itemize}

\begin{tcolorbox}[breakable,colback=teal!4,colframe=teal!35!black,title={Before you move on}]
What does \ja{まちあいしつ} mean? (\textquotedblleft{}A waiting room\textquotedblright{}.) Say it once more.
\end{tcolorbox}
\section[\texorpdfstring{\ja{じどうはんばいき} (jidōhanbaiki)}{じどうはんばいき (jidōhanbaiki)}]{\ja{じどうはんばいき} (jidōhanbaiki) — a vending machine}

\label{lesson:JA-C372-jidouhanbaiki}



\begin{quote}
Before the new one: say the Japanese for to meet up, then the Japanese for a waiting room.
\end{quote}

\subsection*{You'll want to know: \ja{じどうはんばいき}}
\textbf{\ja{じどうはんばいき}} — \emph{jidōhanbaiki} — \textquotedblleft{}a vending machine\textquotedblright{}.

\begin{quote}
\textbf{\ja{じどうはんばいきで} \ja{おちゃを} \ja{かいました。}} — \emph{jidōhanbaiki de ocha o kaimashita} — \textquotedblleft{}I bought tea from the vending machine.\textquotedblright{}
\end{quote}

\begin{cousinweb}
\textbf{\ja{じどうはんばいき}} (\emph{jidōhanbaiki}) is Sino-Japanese (kango): five Chinese characters read with their on'yomi, readings borrowed from Chinese. They are \emph{ji}, \textquotedblleft{}oneself\textquotedblright{} (Middle Chinese \emph{dzijH}), \emph{dō}, \textquotedblleft{}move\textquotedblright{} (\emph{duwngX}), \emph{han}, \textquotedblleft{}trade\textquotedblright{} (\emph{pjonH}), \emph{bai}, \textquotedblleft{}sell\textquotedblright{} (\emph{meaH}) and \emph{ki}, \textquotedblleft{}machine\textquotedblright{} (\emph{kj+j}) (KANJIDIC2; Middle Chinese in Baxter's transcription).
\end{cousinweb}

\begin{grammarlens}[title={What you've built}]
One more word for getting around.
\end{grammarlens}

\subsection*{Your turn}
\begin{itemize}
\raggedright
  \item \emph{Say it:} \emph{jidōhanbaiki}
  \item \emph{Say it:} \emph{jidōhanbaiki}, once more
  \item \emph{Say it:} \emph{jidōhanbaiki de ocha o kaimashita} — I bought tea from the vending machine.
  \item \emph{Recall:} say the Japanese for to meet up, then the Japanese for a waiting room, then say \emph{jidōhanbaiki} again
\end{itemize}

\begin{tcolorbox}[breakable,colback=teal!4,colframe=teal!35!black,title={Before you move on}]
What does \ja{じどうはんばいき} mean? (\textquotedblleft{}A vending machine\textquotedblright{}.) Say it once more.
\end{tcolorbox}
